% 本文件由 LaTeX 排版 Agent 根据有效 translation.json 自动生成。
% author=Hilary of Poitiers; work_id=3303; title=圣咏讲道集

\chapter{圣咏讲道集}

% structure: p0001 source=h1 target=omitted reason=page-title-already-rendered-by-parent

圣咏 1\newline{}圣咏 54/53\newline{}圣咏 131/130\par

\SourceRecord{Translated by E.W. Watson and L. Pullan.；Nicene and Post-Nicene Fathers, Second Series；Vol. 9.；Edited by Philip Schaff and Henry Wace.；Buffalo, NY: Christian Literature Publishing Co.,；1899.；<http://www.newadvent.org/fathers/3303.htm>.}

% structure: p0001 source=h1 target=section reason=page-title

\section{圣咏第一篇讲道}

关于圣希拉略的《圣咏讲道集》，本卷导言（页 xl.-xlv.）已有介绍。对其诠释原则尚待补充数语。此原则可从他本人《圣咏训言》（\WorkTitle{Instructio Psalmorum}）第五节的陈述中得以概括——此乃讲道前的序论：\Quote{毫无疑问，圣咏的语言必须藉福音教导的光照来解释。因此，无论藉谁之口\OriginalTerm{预言之神}{Spirit of prophecy}曾发言，其言语的全部目的，都是为教导我们关于我们主耶稣基督来临的光荣与德能、降生成人的光荣与德能、受难的光荣与德能、王国的光荣与德能以及我们复活的光荣与德能。此外，一切预言对世俗理智和异教智慧都是封闭而封印的，正如依撒意亚所说：\Emph{这一切话语对你们犹如这封住的书卷上的言语}\BibleRef{依 29:11}……整部圣咏集是由寓意和预象的含义交织而成，借此在我们眼前展现了天主独生子的全部奥迹：祂将生于肉身，受难，死去，复活，并永远与那些因相信祂而分享祂光荣的人一同为王，且要审判其余的人类。}诚然，希拉略不时有所区分，有时甚至很敏锐地辨明哪些段落必须如此解释、哪些不可，但大部分注释仍是神学的，因而也是奥秘的。圣咏集并非用于建立教义。希拉略没有其他独立依据的任何立场，都不会仅凭寓意解释来维持，也不从这类寓意中推导任何结论。它们仅仅用于累积性地确认已经以其他方式启示的真理。其结果是一部更多阐发希拉略本人思想而非圣咏作者思想的注释；讲道集虽有伟大优点，却也被明显而严重的缺陷所抵消。当然，它对圣咏写作背景的兴趣甚少。在希拉略眼中，它们本质上是预言，他通常只满足于将作者描述为\Quote{先知}。与历史一样，圣咏集的精神也如此。几乎没有证据表明他认识到其中包含人类对天主最崇高、最完美的虔敬表达，更不用说他对其中诗歌的崇高有所欣赏。对于后一缺失，有充分理由。希腊文《七十贤士译本》和旧拉丁文译本对我们而言具有可敬的古远和神圣的联系，但很难说它们能激发想象力。然而，希拉略固然因主使用过希腊文译本且为普遍公认而视其为权威，但他对待它并非作为文学，而是以律师解释和应用古老宪章条款的精神。拉丁文译本也不太可能感动希拉略，如同它有时感动今日阅读它的我们——我们在其未经修饰的句子中发现某种庄严和力量。其生硬只能使一位能够用他活生生的、受过良好教养的语言清晰思考和表达的人（如我们已在导言 iii. 中所说）感到震惊，其晦涩只能使他困惑。但尽管有种种不利条件，他仍创作了一部伟大而深具基督信仰的作品，具有持久的价值和趣味，并在神学与道德思想上产生了持久的影响。因为在所有这些讲道中——尤其在本卷所译的那些——罗马人在道德反思上的天才得以彰显，并为圣安博在其《论司铎职守》（\WorkTitle{De officiis ministrorum}）这类作品中成功遵循的模式奠定了基础。\par

阅读圣咏的首要知识条件，是能够看出我们该视圣咏作者为谁的代言人，以及他向谁说话。因为它们并非全都具有同一性质，而是出自不同作者、属于不同类型。因为我们常发现天主圣父的位格置于我们面前，如在第八十八篇圣咏中：\BibleQuote{我从百姓中举扬了一位被选者，找着了我的仆人达味，以我的圣油傅抹了他。他必呼叫我：你是我的父亲，是我的救恩的扶持者。我也要立他为长子，高于地上的君王}；而在我们可称之为多数的圣咏中，则引入圣子的位格，如在第十七篇中：\BibleQuote{我所不认识的人民要事奉我}；以及在第二十一篇中：\BibleQuote{他们分了我的外衣，为我的里衣拈阄}。但第一篇圣咏的内容却禁止我们将其理解为父或子的位格：\BibleQuote{但他的旨意在于主的法律，日夜默思祂的法律}。在我们所说指涉圣父位格的那篇圣咏中，所用的措辞完全恰当，例如：\BibleQuote{他必呼叫我：你是我的父亲，我的天主，是我的救恩的扶持者}；而在我们听见圣子说话的那篇中，祂藉着所使用的表达方式自称是这些话的作者，说：\BibleQuote{我所不认识的人民事奉了我}。这就是说，当父说\Quote{他必呼叫我}，而子说\Quote{人民事奉了我}时，他们表明是他们亲口谈论自己。然而在这里，我们看到\BibleQuote{但他的旨意在于主的法律}，显然不是主的位格谈论自己，而是另一人的位格，赞扬那其意志在于主的法律之人的福祉。因此，我们应在此认出先知的位格，圣神藉先知之口说话，通过其嘴唇的工具将我们提升到对属灵奥秘的认识。\par

2. 当他这样说时，我们必须查究他所说的那人是指谁。他说：\BibleQuote{有福的人，不随从不敬虔者的计谋，不站立在罪人的道路，不坐在亵渎者的座位上；他的意愿却在于上主的法律，他昼夜默思他的法律。他必像一棵栽在溪水旁的树，按时结果实，他的叶子也不枯干，凡他所做的无不顺利。}我从与许多人的交谈或他们的书信著作中发现，关于这篇圣咏，许多人的意见认为我们应当理解它为描述我们的主耶稣基督，并且以下经句所赞颂的是他的幸福。但这种解释在方法和推理上都是错误的，尽管无疑它源于一种虔敬的思想倾向，因为全部圣咏集都应指向他：必须凭藉由理性引导的正确知识方法来确定这段经文所指的他生命中的时间和地点。\par

3. 但这篇圣咏开头的言语与圣子的位格和尊严颇不相称，而全部内容本身即是对草率匆忙地用以称扬他的谴责。因为当经上说：\BibleQuote{他的意愿却在于上主的法律}，既然法律是由天主圣子所赐，怎么能把那种因他的意愿在于上主的法律而有的幸福归给那本身就是法律之主的那一位呢？法律属于他，他本人在\BibleBook{}{第七十七篇圣咏}中亲自表明：\BibleQuote{我的百姓，请听我的法律；请侧耳倾听我口中的言语。我要开口说比喻。}福音作者玛窦进一步证实这些话是由圣子所说，他说：\Quote{因此他说话用比喻，为应验那所说的话：\BibleQuote{我要开口用比喻来说话。}}\BibleRef{玛 13:35}上主于是以实际行动实现了自己的预言，用他曾应许说话的那些比喻来讲论。但是，那句话——\BibleQuote{他必像一棵栽在溪水旁的树}——其中以比喻述说幸福的增长，又怎么可能应用于他的位格，并称一棵树比天主圣子更幸福，以及成为他幸福的原因呢？倘若在他与树之间就朝向幸福的增长方面建立了类比，则会出现这种情况。再者，按照\BibleBook{}{箴言}\BibleRef{箴 8:22}和宗徒的教导，他既是在万世之前，也在永恒的时间之前，且是万物的首生者；既然万有都是在他内并藉着他受造的，他又怎么能因变得像他自己所创造的对象而成为幸福的呢？因为创造者的权能并不需要藉与任何受造物的比较而得以提升，首生者的亘古年龄也不容许涉及不合时宜的时间条件的比较，若将他与一棵树相比，便是如此。因为那将在未来某个时刻存在的东西，不能被视为先前已存在或现今存在于任何地方。但已经存在的东西，不需要任何时间延伸才开始存在，因为它从开始直到现在已经拥有连续的存在。\par

4. 因此，既然这些话被认为不适合应用于天主独生子、我们的主耶稣基督的天主性，我们就必须推定，那位在此被先知赞为有福的人，乃是那藉着对正义的热诚和对一切义德的圆满履行，努力使自己与主所取并作为人而生的那身体相似的人。随着圣咏讲解的展开，将表明这是必要的解释。\par

5. 圣神选择这一宏伟而高贵的圣咏集引言，目的是藉着幸福的盼望激发软弱的人对虔敬的纯正热忱，教导他降生成人的天主的奥秘，应许他分享天上的荣耀，宣告审判的惩罚，宣扬双重的复活，显示在他赏报中所见的天主计划。这伟大的预言确实是根据无瑕而成熟的计划奠定基础的；他的旨意是：与有福之人相关的盼望能吸引软弱的人性归向对信德的热忱；树的幸福的类比能成为幸福盼望的保证；他针对不敬虔者所发的义怒的宣告能为不虔敬的过分行为设定恐惧的界限；圣徒集会中地位的差异能标志功绩的差异；为判断义人道路而设立的准则能彰显天主的威严。\par

但我们现在来看主题及表达它的言语。\par

6. \BibleQuote{有福的人，不随从不敬虔者的计谋，不站立在罪人的道路，不坐在亵渎者的座位上；他的意愿却在于上主的法律，他昼夜默思他的法律。}\par

先知列举了五种谨慎，作为经常存在于有福之人思想中的：第一，不随从不敬虔者的计谋；第二，不站立在罪人的道路；第三，不坐在亵渎者的座位上；其次，将他的意愿置于上主的法律中；最后，昼夜默思其中。因此，必须在\InnerQuote{不敬虔者}与\InnerQuote{罪人}之间，在\InnerQuote{罪人}与\InnerQuote{亵渎者}之间作出区分；主要因为此处不敬虔者有计谋，罪人有道路，亵渎者有座位；又因为问题在于\InnerQuote{行走}而非\InnerQuote{站立}于不敬虔者的计谋中；\InnerQuote{站立}而非\InnerQuote{行走}于罪人的道路上。现在如果我们想要理解这些事实的原因，我们必须注意罪人与不敬虔者之间的确切区别，如此才能明白为何罪人被指派了一条道路，而不敬虔者被指派了一个计谋；其次，为何问题在于\InnerQuote{站立}在道路中和\InnerQuote{行走}在计谋中，而人们通常把\InnerQuote{站立}与计谋相连，把\InnerQuote{行走}与道路相连。\par

并非每个罪人都是\OriginalTerm{不孝者}{undutiful}，但不孝者必定是罪人。让我们从普遍经验中取一例。儿子们即使酗酒、放荡、挥霍，仍能爱他们的父亲；虽有这一切恶习，因此不免罪责，却仍可能不犯不孝之罪。但不孝者，即使可能是节制和节俭的典范，仅因藐视父母这一点，就比犯下不孝之外所有罪过的人更为恶劣。\par

7. 毫无疑问，正如这个例子所证明的，\OriginalTerm{不敬者}{ungodly}（或不虔者）必须与罪人区分开来。事实上，普遍意见一致称那些不屑寻求认识\Person{天主}的人为不敬者，他们怀着不恭之心想当然地以为没有世界的造物主，断言世界是通过偶然运动达到我们所见的秩序与美，为了剥夺其造物主对善生或罪生的一切审判权，他们主张人仅仅通过自然法则的作用而进入存在并离开它。\par

因此，所有这些人的谋虑是摇摆不定、不稳妥、含糊不清的，在同样熟悉的道路上和同样熟悉的地面上徘徊，从未找到安息之处，因为它未能达成任何确定的决定。他们从未在自己的体系中上升到世界造物主的教义，因为他们非但没有回答我们关于世界的原因、开端和持续，世界是为着人还是人为着世界，死亡的原因、范围与本质等问题，反而在这\OriginalTerm{无神论争}{godless argument}的圆圈内不停地转动，在这些幻想中找不到安息。\par

8. 此外，还有不敬者的其他谋略，即那些陷入异端者，不受\WorkTitle{新约}或\WorkTitle{旧约}律法的约束。他们的推理总是走恶性循环的路线；没有抓握或立足之点来止住自己，他们踏着无休止的循环，无尽地优柔寡断。他们的不敬在于衡量\Person{天主}，不是按祂自己的启示，而是按他们自己选择的标准；他们忘记，制造一位\Person{天主}和否认祂一样是不敬；如果你问他们这些意见对信德和望德有何影响，他们便困惑混乱，离题游荡，故意回避辩论的真正问题。因此，那不曾行走在这类不敬者的谋略中的人，不，甚至不曾有过行走其中的愿望的人，是有福的，因为哪怕一刻思想不敬之事也是罪过。\par

9. 下一个条件是，那不曾行走在不敬者谋略中的人，不应站在罪人的道路上。因为有许多人，他们对\Person{天主}的宣认虽然宣告他们不犯不敬之罪，却不能使他们脱离罪；例如那些留在教会内却不遵守其律法的人；如贪婪者、酗酒者、争吵者、放荡者、骄傲者、伪善者、说谎者、掠夺者。无疑，我们因自然本能的冲动而被推向这些罪；但对我们有益的是，从被催逼进入的那条路上退出，不在其中站立，因为我们被提供了如此容易的逃脱之路。正是为此，那没有站在罪人道路上的人是有福的，因为自然将他推入那条路，而宗教信仰将他拉回。\par

10. 现在，获得幸福的第三个条件是：不坐在朽坏的座位上。法利塞人坐在\Person{梅瑟}的座位上教导，\Person{比拉多}坐在审判的座位上：那么，我们要认为什么座位是朽坏的？当然不是\Person{梅瑟}的座位，因为主指责的是座位上的占据者，而非座位的占据，祂说：\BibleQuote{经师和法利塞人坐在\Person{梅瑟}的座位上；凡他们对你们所说的，你们要遵行；但不要效法他们的行为。}\BibleRef{玛 23:2} 那个座位的占据不是朽坏的，因为主亲自命令要顺从它。那实在必须是朽坏的，\Person{比拉多}洗手试图避免其感染。因为许多人，甚至敬畏\Person{天主}的人，因追求世俗荣誉而被误导；他们渴望管理法庭的法律，尽管他们受教会法律的约束。\par

但尽管他们履行职责时带有宗教意向，正如他们仁慈而正直的态度所表明的，他们仍然无法逃脱一种来自其一生所从事事务的传染性感染。因为民事案件的处理不容许他们忠于教会法律的神圣原则，即使他们愿意如此。并且，在不放弃他们虔诚目的的情况下，他们被自己所获得的座位之必要条件所迫，违背意愿地时而使用斥责，时而侮辱，时而惩罚；正是他们的位置使他们既成为约束他们必然性的创造者，也成为其受害者，他们的体系仿佛浸透了感染。因此，先知称他们的座位为\Emph{朽坏的座位}，因为它的感染毒害了宗教心者的意志。\par

11. 但他没有随从不虔敬者的计谋，没有站在罪人的道路上，没有坐在疫病者的座位上\BibleRef{咏 1:1}，这一事实并不构成那人幸福的圆满。因为相信独一天主是世界的创造者，通过追求谦逊的善行而避免罪恶，偏爱隐居生活的宁静胜过公共地位的显赫——这一切甚至在外邦人身上也可以找到。但在这里，先知在描绘按天主肖像而造的完美之人时——一个可以作为永恒幸福之崇高榜样的人——指出他所实践的不是寻常的美德，并指出为了获得圆满的幸福，\BibleQuote{\Emph{但他的意愿是在上主的法律中}}\BibleRef{咏 1:2}。戒除先前所提到的那些事是无用的，除非他的心思集中于下文所说的\BibleQuote{\Emph{但他的意愿是在上主的法律中}}。先知不寻求恐惧。大多数人因恐惧而遵守法律的界限；少数人因意愿而服从法律：因为恐惧的特征是不敢省略它所害怕的事，而完美的虔诚的特征是乐意服从命令。这就是为什么那人的意愿（而非恐惧）是在天主的法律中的人是有福的。\par

12. 但有时意愿需要补充；仅仅渴望圆满的幸福并不能赢得它，除非实行紧随意向之后。圣咏接着说：\BibleQuote{\Emph{他将在祂的法律中昼夜默想}}。那人通过持续不懈的默想法律而达成幸福的圆满。现在可能有人反对说，由于人性软弱的条件（需要休息、睡眠和饮食的时间），这是不可能的：因此我们身体的状况使我们无法获得幸福的希望，因为身体的需用中断了我们的日夜默想。与此段相似的是宗徒的话：\BibleQuote{\Emph{你们要不断祈祷}}\BibleRef{得前 5:17}，好像我们被要求无视身体的需要而无间断地继续祈祷！因此，在法律中的默想不在于念诵其词句，而在于虔诚地遵行其诫命；不在于仅仅阅读书籍和文献，而在于对其内容的实践默想和操练，并通过我们日夜所做的行为来满全法律，正如宗徒所说：\BibleQuote{\Emph{所以你们或吃或喝，或无论做什么，一切都要为光荣天主而做}}\BibleRef{格前 10:31}。确保不断祈祷的方式是：每个虔诚的人通过在天主眼中悦纳且始终为祂的光荣而做的工作，使自己的生命成为一种漫长的祈祷；因此，昼夜按照法律生活的本身将成为昼夜在法律中的默想。\par

13. 但现在，那人由于远离不虔敬者的计谋、罪人的道路和疫病者的座位，并由于甘心乐意昼夜默想天主的法律而获得了圆满的幸福，接下来要向我们展示他所赢得的幸福将结出丰硕的果实。因为幸福的预尝包含着未来幸福的胚芽。因为下一节经文说：\BibleQuote{\Emph{他像一棵栽在溪水旁的树，按时结果，叶子也不枯落}}。这也许会被视为荒谬且不恰当的比喻，其中赞美了一棵栽种的树、溪水、结果实、自己的季节和不枯落的叶子。这些在世俗的判断看来可能微不足道。但让我们审视先知的教导，看看用来描绘幸福的事物和话语中所蕴含的美。\par

14. 在\BibleBook{}{创世记}\BibleRef{创 2:9}中，当立法者描绘天主所栽种的乐园时，显示每棵树都悦人眼目且好作食物；也提到在园子中间有一棵生命树和一棵知善恶树；接着乐园被一条河流灌溉，它后来分成四道。所罗门先知在关于智慧的劝勉中教导我们什么是生命树：\BibleQuote{\Emph{她为所有抓住她、倚靠她的人是一棵生命树}}\BibleRef{箴 3:18}。这棵树因此是活的；不但活着，而且更进一步，受理性指导；受理性指导，即就结果子而言，并且不是偶然或不当时节，而是按时。这棵树栽在天主王国领域内的溪水旁，当然就是在乐园里，在河流流出后分成四道的地方。因为他说不是\Emph{在溪水后面}，而是\Emph{在溪水旁}，即在四道河流最初各自接受水流的地方。这棵树栽在这样一个地方，主——祂就是智慧——带领承认祂为主的强盗去那里，说：\BibleQuote{\Emph{我实在告诉你，今天你就要同我在乐园里}}\BibleRef{路 23:43}。现在我们已经依据先知权威表明，智慧即是基督，按照即将来临的降生成人和苦难的奥迹，被称为生命树，我们必须从福音中进一步为这一解释的严格真实性寻找支持。当犹太人说他藉贝尔则布驱魔时，主亲口将自己比作一棵树：\BibleQuote{\Emph{你们要么使树好，它的果子也好；要么使树坏，它的果子也坏，因为树凭果子就可以认出}}\BibleRef{玛 12:33}；因为虽然驱魔是一个好果子，他们却说他是贝尔则布，他的果子是可憎的。他也毫不迟疑地教导说使树幸福的能力存在于祂本人内，当祂前往十字架的路上说：\BibleQuote{\Emph{如果他们对待青绿的树木尚且这样，那枯槁的将会怎样呢？}}\BibleRef{路 23:31}，藉着这青绿树的形象宣告在祂里面没有任何东西会受死亡的枯槁所支配。\par

15. 那位有福的人将变得像这棵树，当他被移植，如同那盗贼一样，进入园中，栽在水渠旁：他的种植将是那种幸福的、不能被拔除的新种植，主在福音中诅咒另一种种植时曾提到它，说：\BibleQuote{凡不是我父所种植的，必要被拔除。}\BibleRef{玛 15:13}因此，这棵树将结出它的果实。如今，在圣经中天主圣言从树的果实教导教训的所有其他地方，都提到它们制造果实而非结出果实，如它所说：\Quote{好树不能制造坏的果实}，以及在依撒意亚书中对葡萄园的抱怨：\BibleQuote{我指望它结葡萄，它反倒长了荆棘。}\BibleRef{依 5:2}但这棵树将结出它的果实，因为它为此被赋予了自由意志和理解力。因为它将在自己的季节结出果实。请问，在什么季节？当然是在宗徒所说的季节：\BibleQuote{使你们知道祂旨意的奥秘，是按照祂在祂自己所预定的美意，在时期圆满的经纶中。}\BibleRef{弗 1:9}那么，这就是时期的经纶，它调节领受者领受的正当时刻和施予者施予的正当时刻；因为施予者选择季节。但时间上的延迟取决于时期的圆满。因为结出果实的经纶等候时期的圆满。那么，你要问，这要分施的果实是什么？无疑就是同一位宗徒所说的：\BibleQuote{祂将改变我们卑贱的身体，使之相似祂光荣的身体。}\BibleRef{斐 3:21}这样，祂将把祂的果实赐给我们，这些果实祂已经在那位祂所拣选的人（就是以树为形象所描绘的那人）身上实现了完美，祂已经完全除去了他的必朽性，并使他分享祂自己的不朽。\par

那时，这个人将像那棵树一样有福，当他最终站立在天主的荣耀中，被塑造得相似于主。\par

16. \Quote{但这棵树的叶子必不落下。}这并不奇怪，因为它的果实也不会落地，要么因为成熟而被迫落下，要么因外力而被摇落，而是它将结出它们，通过一种理性的服务来分配它们。现在，叶子的属灵意义通过基于物质对象的比较得以阐明。我们看到，叶子是为了保护果实而发芽簇拥在其周围，目的是保护和形成一道围栏给幼嫩的枝条。那么，叶子所象征的就是天主言语的教导，在其中应许的果实被包裹。正是这些言语仁慈地遮蔽我们的希望，保护它们免受这世界暴风的侵袭。这些叶子，即天主的言语，必不落下，因为主亲自说过：\BibleQuote{天地要过去，但我的话绝不会过去。}\BibleRef{玛 24:35}因为天主所说的话，没有一句会落空或落下。\par

17. 如今，我们所谈论的树的叶子并非毫无价值，而是给万民带来健康的泉源，这由圣若望在\WorkTitle{默示录}中作证，他说：\BibleQuote{祂又给我显示了一条生命水的河流，明亮如水晶，从天主和羔羊的宝座中流出来；在街道的中央和河的两岸，有生命树，结十二样果子，每月结出它的果实；树上的叶子可医治万民。}\BibleRef{默 22:1}\par

物质的表现如此揭示天上的奥秘，尽管物质本身不能传达全部属灵意义，但只从物质方面看待它们就是损害它们。我们本应期望听到，在显示给圣人的河的两岸，有许多树，而不是一棵树。但因为生命树在圣洗圣事中总是在每种情况下都是一棵，为从各方面来到它那里的人提供宗徒信息的果实，所以在河的两岸只站着一棵生命树。在天主宝座中间看到一只羔羊，一条河，一棵生命树：三个形象中包含了降生成人、圣洗圣事和受难的奥秘，它们的叶子——即福音的言语——通过不能落空的教导给万民带来医治。\par

18. \Quote{凡他所做的，尽都顺利。}从此以后，祂的恩赐和法令再不会被藐视，如同在亚当身上那样，他因违犯法律而犯罪，失去了确定不朽的幸福；但现在，由于生命树的救赎，即主的受难，我们身上发生的一切都是永恒的，并因我们将来相似生命树而永恒地意识到幸福。因为他们所做的一切都将顺利，不再在变化和人的软弱中进行，因为腐朽将被不朽吞没，软弱被永生吞没，属土的形状被天主的形状吞没。那么，那棵栽种并按时结果子的树，将有福的人与之相似，他自己也被栽种在园中，使天主所栽种的得以存留，永不被拔除，在园中天主所做的一切都将被导向顺利的结果，远离属于人的软弱和时间的腐朽，那必须被拔除的。\par

19. 先知在陈述了那人圆满的幸福之后，接着便要宣告那留给不虔敬者的惩罚是什么。因此接续说：\Quote{不虔敬的人却不是这样，而是像糠秕，被风吹散}。不虔敬的人绝无希望获得那幸福树的形象；他们惟一的命运就是飘荡、簸扬、压碎、分散和不安；从他们身体状态的坚固框架中被抖落出来，必被扫到惩罚中，成为尘土，成为风的玩物。他们不会归于虚无，因为惩罚必须在他们身上找到一些可作用的东西，但被磨成微粒、轻飘、虚幻、干燥，他们将不断被抛来抛去，成为永不休息的惩罚的玩物。他们的惩罚由同一位先知在另一处记载，他说：\Quote{我要将他们打碎，如同风前的尘土，如同街上的泥土，我要消灭他们}。\par

因此，正如幸福有指定的样式，惩罚也有其样式。因为风吹散尘土并非难事，人们走过街上的泥泞也几乎不觉踩过它；同样，地狱的惩罚容易毁灭和驱散不虔敬者，他们罪恶的逻辑结果就是将他们化为泥浆，压成尘土，失去一切坚固的实质，因为他们本就是尘土和泥浆，且仅仅是尘土和泥浆，除了受罚之外一无用处。\par

20. 先知看到他们坚固的实质化为尘土，将使他们完全失去分享那树按季节赐给幸福之人的果实，因此接着加上：\Quote{因此，不虔敬的人在审判中不得复活}。他们不得复活，并非宣告这些人归于无有，因为他们确实将作为尘土存在；被拒绝的是复活受审。不存在并不能使他们免于惩罚的痛苦；因为那不存在的会逃脱惩罚，而他们却存在以受罚，因为他们将成为尘土。成为尘土，无论是被晒成尘土还是被磨成尘土，都不是失去存在状态，而是状态的改变。但他们不得复活受审这一点表明，他们失去的不是复活的能力，而是复活受审的特权。至于复活受审的特权究竟指什么，主在福音中宣告说：\Quote{信我的人，不受审判；不信的人，已受了审判。这审判就是：光明来到了世界，世人却爱黑暗甚于光明} \BibleRef{若 3:18}。\par

21. 主这番话的措辞使漫不经心的听众和粗心草率的读者感到困惑。因为祂说：\Quote{信我的人，不受审判}，将信徒排除在审判之外；又说：\Quote{不信的人，已受了审判}，将不信者排除在审判之外。如果祂已经这样使信徒免于审判、将不信者排除在外，既不让这一类人也不让那一类人有受审的机会，那么当祂接着加上第三句：\Quote{这审判就是：光明来到了世界，世人却爱黑暗甚于光明}时，又怎能被认为是前后一致的呢？因为显然没有留下审判的余地，既然信徒和不信者都不受审判。这无疑会是漫不经心的听众和草率的读者得出的结论。然而，这番话自有其恰当的意义和合理的解释。\par

22. 基督说，信的人不受审判。那么审判一个信徒有何必要呢？审判出于模棱两可；当模棱两可消失时，就没有必要进行审讯和审判了。因此，甚至不信者也不需要受审，因为他们是毫无疑义的不信者；但在同样使信徒和不信者免于审判之后，主又加上了需要审判的情况和人。因为有些人处于敬虔和不虔敬之间，与两者都有相似之处，但又不严格属于任何一类，因为他们是由两者的混合所构成的。他们不能被归入信仰的行列，因为他们心中有某种不信的掺杂；他们也不能被列于不信之中，因为他们并非没有一部分信仰。许多人因敬畏天主而留在教会的围墙之内，但同时却因世界的诱惑而被引诱犯世俗的过错。他们祈祷，因为害怕；他们犯罪，因为出于本意。对来生美好的盼望使他们自称为基督徒；对现世欢乐的诱惑使他们行事如外邦人。他们不滞留于不虔敬，因为他们尊敬天主的名；他们又不虔敬，因为他们追求与虔敬相悖的事。他们不能不最爱那些永远不能使他们名副其实的事物，因为他们行这类事的渴望强于忠于自己名字的渴望。因此，主在说了信徒不受审判、不信者已受审判之后，又加上：\Quote{这审判就是：光明来到了世界，世人却爱黑暗甚于光明}。\par

这些就是那些等候审判的人，这个审判是不信者早已受过的、信徒所不需要的：因为他们爱黑暗甚于光明；并非他们不爱光明，而是因为他们对黑暗的爱更为强烈。因为当两种爱彼此竞争时，总有一种得到偏爱；他们的审判源于这样一个事实：虽然他们爱基督，却更爱黑暗。这些人都将受审；他们不像敬虔者那样免于审判，也不像不虔敬者那样已受审判；审判正等候他们，因为他们蓄意选择了所爱的。\par

23. 先知所遵循的正是福音中所规定的这个计划和秩序，当时他说：\BibleQuote{因此，不虔敬的人在审判中不得复活，罪人也不得进入义人的集会。}他在不虔敬的人身上不留审判，因为他们已被审判过了；另一方面，他不准罪人——我们在前一篇讲道中已说明他们与不虔敬的人有别——得享义人的集会，因为他们将要受审判。因为不虔敬使前者预先被审判，而罪却使后者留在日后受审判。如此，不虔敬者既已被审判，就不被接受进入对罪人的审判；而罪人既将受审判，就被认为不配享有那些不会再受审判的义人的集会。\par

24. 这种区分的根源在于以下的话：\BibleQuote{因为主认识义人的道路，但不虔敬的人的道路必要灭亡。}罪人不能接近义人的集会，原因是主认识义人的道路。祂认识，并不是从无知进到知识，而是因为祂愿意屈尊认识。因为天主并没有人类情感的波动，时而认识、时而不认识任何事物。蒙福的宗徒保禄指出我们是如何被天主认识的，他说：\BibleQuote{若你们中有人是先知或属神的人，就该承认我写给你们的，是主的诫命；但若有人不承认，他也不被承认。}\BibleRef{格前 14:37}\par

如此他表明，那些认识天主之事的人是被天主认识的：当他们认识时，他们就要被认识，也就是说，当他们借着已知的虔敬之功劳，达到被认识的尊荣时，好使这认识被视为被认识者方面的成长，而非不认识者方面的成长。\par

现在天主藉着亚当今亚巴郎的实例清楚表明：祂不认识罪人，却认识信徒。因为当亚当犯罪时，有人对他说：\BibleQuote{亚当，你在哪里？}\BibleRef{创 3:9}并非因为天主不知道那人仍在园中，而是藉着询问他在哪里，来表明他因犯罪而不配为天主所认识。但亚巴郎长期不为人知——天主的圣言在他七十岁时临到他——因他对主显出了忠信，便藉着以下至高的俯就而得以与天主亲近：\Quote{现在我知道你敬畏天主，为了我，你没有怜惜你的独生子。}\par

天主当然并非不晓得亚巴郎的信德，在亚巴郎因相信依撒格的诞生而信德已被算为他的义德时，祂已知道；但现在，因为他献子时给了敬畏的一个显著证据，他终于被认识、被认可，成为配得不再不为人知。因此，天主也这样认识和不认识——罪人亚当不被认识，而忠信的亚巴郎被认识，也就是说，他配得被那确知万事的天主所认识。因此，义人——他们将不受审判——的道路被天主认识；这就是为什么那些将要受审判的罪人远离他们的集会；而不虔敬的人不会复活起来接受审判，因为他们的道路已经灭亡，并且他们已经由那一位审判了，祂说：\Quote{父不审判任何人，却把审判全交给了子}，我们的主耶稣基督，祂是永受赞美的，直到永远。阿们。\par

\SourceRecord{Translated by E.W. Watson and L. Pullan.；Nicene and Post-Nicene Fathers, Second Series；Vol. 9.；Edited by Philip Schaff and Henry Wace.；Buffalo, NY: Christian Literature Publishing Co.,；1899.；<http://www.newadvent.org/fathers/3303001.htm>.}

% structure: p0001 source=h1 target=section reason=page-title

\section{圣咏第五十三篇}

\Emph{交与歌咏长，达味的训诲诗，当齐弗人来对撒乌耳说：\InnerQuote{达味岂不是在我们中间藏身吗？}}\par

\Emph{天主，求以祢的名拯救我，以祢的权能为我伸冤。天主，求祢俯听我的祈祷，侧耳听我口中的言词。}等等。\par

1. 福音的教理为圣洁有福的达味先知所熟知，他虽然肉身生活在法律之下，却尽己所能地满全了宗徒般的训诲，并证实了天主为他所作的见证：\BibleQuote{我找到了合我心意的人，叶瑟的儿子达味。}\newline{}他不以战争报复敌人，也不以武力对抗谋害他的人，而是效法主的榜样——主的名和祂的温良他都预表了——当他被出卖时，他恳求；当他遇险时，他歌唱圣咏；当他遭人怨恨时，他喜乐；因此他被寻得为合天主心意的人。因为尽管十二营的天使可以在主受难时前来帮助，但为了完美地完成谦卑服事的使命，祂将自己交于痛苦和软弱，只以祈祷说：\BibleQuote{父啊，我把我的灵魂交在祢手中。}\BibleRef{路 23:46} 达味同样如此，他的实际苦难先知性地预言了主将来的苦难，他不以言语或行为对抗敌人；为遵行福音的命令，他不以恶报恶，在痛苦、被出卖、争战中，他效法师傅的温良，呼求主，只以主的武器与不敬虔之人争斗。\par

2. 此篇圣咏的前言引用了一个历史事件；但在描述事件之前，我们受到了关于其范围、时间和所蕴含事件之应用的教导。首先我们有：\Quote{交与歌咏长，达味的训诲诗}。然后接着：\Quote{当齐弗人来对撒乌耳说：\InnerQuote{达味岂不是在我们中间藏身吗？}}这样，达味被齐弗人出卖等待其最终解释。这表明实际发生在达味身上的事包含了对未来之事的预像；一个无辜的人受辱骂，先知被嘲弄的话语讥讽，受天主认可的人被要求处死，君王被出卖给敌人。同样，主被祂本来应当安全的人出卖给黑落德和比拉多。因此，此篇圣咏等待最终解释，并在真正的达味——法律之终——身上找到其意义，那达味拿着钥匙，打开知识之门，实现达味所预言的一切。\par

3. 按希伯来文的准确意义，这个专名的含义为我们解释这段经文提供了不小的帮助。\OriginalTerm{齐弗人}{Ziphims}意思是\Quote{脸上的洒水}，在希伯来文中称为\Quote{齐弗人}。按照法律，洒水是洗净罪恶；它通过信德用洒血来洁净百姓，正如同一位有福的达味所说：\BibleQuote{祢用牛膝草洒我，我就洁净了。}法律通过信德，以全燔祭的血暂时提供了将来主血洒净的预像。但这百姓，像齐弗人一样，只在脸上受洒而不在信德中，只在嘴唇上接受洁净的水珠而不在心里，便对自己的达味变得不信和出卖，正如天主借先知所预言的：\BibleQuote{这百姓用嘴唇尊敬我，他们的心却远离我。}他们准备出卖达味，因为他们心里的信德已死，他们用虚伪的脸履行了法律中所有神秘的仪式。\par

4. \BibleQuote{天主，求以祢的名拯救我，以祢的权能为我伸冤。天主，求祢俯听我的祈祷，侧耳听我口中的言词。}\par

根据我们对标题的解释，先知达味的受苦是我们天主和主耶稣基督受难的一个预像。这就是为什么他的祈祷在意义上与那位成为血肉的圣言的祈祷相对应：祂按人的方式承受一切，按人的方式说话；祂担当了人的软弱和罪恶，以属于人的谦卑向天主祈祷。即使我们不愿意且迟钝接受，这种解释也被话语的意义和力量所要求，以至于毫无疑问，整篇圣咏中的一切都是由达味作为祂的代言人所说。因为他说：\BibleQuote{天主，求以祢的名拯救我。}天主独生子，同时要求祂在万世以前拥有的荣耀，在身体卑微中祈祷，用祂自己先知的话。祂求以天主的圣名拯救祂，这名是祂被称呼且在其中所生的，以便那原本属于祂先前本性和类型的圣名能够拯救祂在祂所出生的身体中。\par

5. 由于整段话是那一位以仆人形像发出的言语——那仆人顺服至死于十字架——祂担负了这形像，并为此恳求属于天主的圣名之拯救助佑，且因那圣名确信得救，便立即加上：\Quote{\Emph{并凭祢的能力审判我}}。因为如今作为祂空虚自己、取了仆人形像之谦卑的酬报，祂在同样谦卑中（祂曾以此谦卑取了这形像）请求恢复祂与天主共享的形像，并已拯救那承受天主圣名的人性——祂曾以天主的身份屈尊降生其中，且是顺服的。为教导我们这圣名（祂藉此祈祷获得拯救）的尊贵并非空洞头衔，祂祈祷要受天主能力的审判。因为公义的判决是审判的必然结果，如圣经上说：\BibleQuote{\Emph{成为顺服至死}\Emph{，且死在十字架上。因此，天主极其举扬祂，赐给祂那超越一切名号的圣名}}。因此，首先那超越一切名号的圣名赐给了祂；其次，这是一个具有决定性力量的审判，因为凭天主的能力，祂——原为天主、却以人而死后——作为人从死亡中复活成为天主，如宗徒所说：\BibleQuote{\Emph{祂由于软弱被钉死，却由于天主的能力而生活}}\BibleRef{格后 13:4}，又说：\BibleQuote{\Emph{因为我不以福音为耻；这福音正是天主的能力，为使一切相信的人获得救恩}}\BibleRef{罗 1:16}。因为藉着审判的能力，人性软弱被拯救以承受天主的圣名和本性；如此，作为祂顺服的酬报，祂被这审判的能力举扬至天主圣名的拯救保护之下；因此祂同时拥有天主的圣名和能力。再者，如果先知是以常人说话的方式开始这段话，他会请求按仁慈或恩惠受审判，而非按能力。但那一位——祂作为天主子由童贞女所生成为人子，如今作为人子，要藉审判的能力恢复天主子圣名和能力——则必须受能力审判。\par

6. 接下来接着：\Quote{\Emph{天主，求祢垂听我的祈祷，侧耳听我口中的言语}}。先知按常理应当说：\Quote{\Emph{天主，求祢垂听我}}。但既然他以那唯一知道如何祈祷者的喉舌说话，我们便看到不断重复的要求：祈求被垂听。圣保禄的话教导我们，没有人知道应当怎样祈祷：\BibleQuote{\Emph{因为我们不知道如何按应当的方式祈祷}}。因此，人在软弱中无权要求他的祈祷被垂听：因为连外邦人的导师也不知道祈祷的真正目的和范畴，且是在主赐予了范本之后。这里所显示的是那一位的完全信心——祂唯独看见圣父，唯独认识圣父，唯独能彻夜祈祷——福音告诉我们主整夜在祈祷——祂在言语的镜中，以我们祈祷时使用的简单言语，向我们展示了最深奥奥秘的真实形像。因此，在提出祂的祈祷应当被垂听的要求时，祂加上了（为教导我们这是祂完全信心的特权）：\Quote{\Emph{侧耳听我口中的言语}}。如今，有谁能认为这是人的信心，竟能这样渴望他口中的言语被垂听？例如那些言语，我们藉以表达心灵的冲动与本能：或愤怒燃烧我们，或仇恨驱使我们毁谤，或痛苦驱使抱怨，或谄媚使我们奉承，或者得利的希望或真理的羞耻产生谎言，或委屈的怨恨产生侮辱？难道曾有任何人一生中完全纯洁忍耐，以至于不会陷入这些人性不稳的软弱？唯独那一位能满有信心地渴望这一点——祂没有犯罪，口中没有诡诈，将背转给击打者，不转脸躲避掌掴，不怨恨嘲笑和唾弃，从不违背祂的旨意——祂完全顺服那安排万事的旨意。\par

7. 祂接着加上为何祈求祂的言语被垂听：\Quote{\Emph{因为陌生人起来攻击我，强暴者寻索我的性命；他们没有将天主放在眼前}}。天主的独生子，天主的圣言，天主且是圣言——虽然无疑祂自己能做圣父能的一切，如祂说：\BibleQuote{\Emph{凡圣父所做的，圣子也同样做}}\BibleRef{若 5:19}，而那描述祂神圣本性的圣名不可分割地含蕴着对神圣能力的不可分离持有——但为给我们呈现人性谦卑的完美榜样，祂既祈祷又承受了人所命定的一切。分享我们的共同软弱，祂祈求圣父拯救祂，为教导我们：祂作为人出生在人性软弱的一切条件下。这也就是为何祂饿、渴、睡觉、疲倦、避开恶人的集会，为何祂忧愁、哭泣、受苦并死亡。为使祂承受这一切条件——不是出于本性，而是出于摄取——在承受这一切后祂复活了。因此，圣咏中一切怨叹都出自属于我们本性的心理状态。如果我们这样理解圣咏的言语，也不必惊奇，因为主亲自见证（如果我们相信福音），圣咏曾以预言方式预告了祂的苦难。\par

8. 如今他们是\Quote{兴起攻击祂的陌生人}。因这些人不是亚巴郎的子孙，也不是天主的子女，而是毒蛇的种类，罪恶的奴仆，客纳罕的后裔，他们的父亲是阿摩黎人，母亲是赫特人的女儿，从他们的祖先魔鬼那里继承了魔鬼般的欲望。此外，强暴者追寻祂的灵魂，如黑落德问司祭长基督应生在何处，又如整个会堂作假见证反对祂。但因他们以为这灵魂属于人性软弱，便\Quote{不把天主放在眼前}；因为天主从祂作为天主时所处的地位，俯就至人类出生的开端；也就是说，祂成为\OriginalTerm{人子}{Son of Man}，而祂原是\OriginalTerm{天主子}{Son of God}，\OriginalTerm{天主的形体}{the Form of God}脱去它原有的形态，成为它原先所不是的，为能进入属于它自己的灵魂和身体。因此，祂既是天主子，也是人子；既是天主，也是人。换言之，天主子带着人类出生的属性而诞生，\OriginalTerm{天主的本性}{the Nature of God}屈尊去接受一个完全由灵魂和肉体所塑造的人的属性。所以，当陌生人兴起攻击祂，当强暴者追寻祂的灵魂时（这灵魂在福音中常是忧伤和沮丧的），他们不把天主放在眼前，因为正是那永恒存在的天主和天主子，带着人类本性的属性而诞生，即带着我们的身体和我们的灵魂，藉由\OriginalTerm{童贞生育}{virgin birth}而成为人。祂所行的伟大光荣的奇迹，从未开启他们的眼睛，使他们看见他们所追寻灵魂的人子，在作为天主子而永恒存在之后，又作为人而有了生命的开始。\par

9. 休止符的引入标志着说话者的改变。祂不再说话，而是被称呼。因为现在预言的话语具有普遍性。因此，在向天主的祈祷之后，祂随即加上，为表明说话者的信心就在祈求的当下已获得了所求的：\Quote{看，天主是我的助佑，主是扶持我灵魂的。祂已向我的仇敌报应了邪恶。} 祂为每一单独的恳求分派了适当的结果，以此教导我们：天主不会忽视聆听，而期待祂怜悯的保证——即祂垂听我们各个恳求——并非不合理之事。因为对于\Quote{陌生人已起来攻击我}这句话，相应的陈述是：\Quote{天主是我的助佑；}而对于\Quote{强暴者已追寻我的灵魂}，祂祈祷蒙垂听的确切结果用以下话语表达：\Quote{主是扶持我灵魂的；}最后，\Quote{他们未将天主放在眼前}这句话，适当地与\Quote{祂已向我的仇敌报应了邪恶}相平衡。因此，天主既帮助抵御那些起来攻击的人，又在强暴者追寻祂圣者的灵魂时扶持它，并且当祂不被放在眼前、不被不敬虔的人考虑时，祂将仇敌所行的恶报应在他们身上。如此，当他们不思虑天主而追寻义人的灵魂并起来攻击祂时，祂却得救并受扶持，而他们却发现，那被他们专注于邪恶行为时所不考虑的祂，却将他们的恶意转向他们自己而施行报复。\par

10. 因此，让纯正的宗教怀有这种信心，不要怀疑：在来自人的迫害和灵魂的危险中，它仍然有天主作为它的助佑，知道如果最终它遭遇暴力和不义的死亡，灵魂离开身体的帐幕时，必在天主——它的扶持者——那里安息；此外，让它对报应有完全的把握，想到一切恶行都回到作恶者的头上。天主不能被指控为不义，完美的善不受邪恶意志的冲动和动静所玷污。祂并不是出于恶意而唤醒灾祸，而是以报复来报应；祂降下灾祸并非因祂希望我们受苦，而是针对我们的罪。因为这些邪恶普遍被设立为惩罚的工具，却不毁灭生命，这是公义审判的严厉公正的规条。但这些邪恶因正义的律法而从义人身上被挡回，并因那审判的正义而转向不义之人。每一过程都同样公正：对于义人，因为他们为义，便只是警告性地展示邪恶而不实际降下；对于恶人，因为他们罪有应得，便加以惩罚性的邪恶；义人虽被展示，却不会遭受它；恶人则永不停止遭受它，因为它被展示给他们。\par

11. 此后，又回到最初向之提出祈求的天主位格：\Quote{求祢以祢的真理消灭他们}。真理驳斥谎言，谎言被真理消灭。我们已表明，前述全部祈祷乃是那人类本性——天主子由此本性而生——的发声；因此，这里是人性的声音，呼求\OriginalTerm{天主圣父}{God the Father}，以祂的真理消灭祂的仇敌。这真理是什么，无可置疑；它当然是那位曾说的：\Quote{我是生命、道路、真理}。\BibleRef{若 14:6} 仇敌被真理消灭，因为当他们试图以假见证赢得基督的定罪时，却听见祂已从死者中复活，不得不承认祂以\OriginalTerm{天主性}{Godhead}的全部真实恢复了祂的荣耀。不久，他们因饥荒和战争而遭毁灭，找到了他们钉死天主的回报；因为他们将生命之主判处死亡，并毫不在乎在祂身上通过祂的荣耀工作所彰显的天主真理。因此，当天主真理在祂复活、恢复祂圣父光荣的威严，并证明祂所拥有的完美天主性的真实时，便消灭了他们。\par

12. 鉴于我们反复且一贯地断言，被高举在十字架上的乃是天主的独生子，被判死刑的那位因祂由永恒圣父所生的本性之起源而是永恒的，必须清楚理解：祂受苦并非出于本性的必然，而是为完成人类救恩的奥迹；祂是自愿而非被迫受苦。虽然这苦难不属于祂作为永恒圣子的本性——天主的不变性使其免于任何有损尊严的侵袭——但它是自由承担的，旨在实现惩罚的职能，然而并不使受苦者感受到惩罚的痛苦：并非这苦难不引起痛苦，而是因为天主性感觉不到痛苦。因此，天主是透过自愿忍受而受苦；但尽管祂完全经历了苦难的全部力量（这力量必然使受苦者感到痛苦），祂却从未放弃其本性的能力以致感受到痛苦。\par

13. 接着经文说：\BibleQuote{我要向祢自由地献祭。} 法律中的祭献，包括全燔祭和山羊、公牛的祭品，并不涉及自由意志的表达，因为凡违反法律的都被判处诅咒。谁不献祭，就使自己受诅咒。而且必须始终完成全部的祭献行动，因为诫命中加上诅咒，禁止人在献祭的义务上懈怠。我们的主耶稣基督正是从这诅咒中救赎了我们，正如宗徒所说：\BibleQuote{基督救赎了我们脱离法律的诅咒，为我们成了诅咒，因为经上记载：凡挂在木架上的，是可诅咒的。}\BibleRef{迦 3:13} 因此，祂将自己交付于被诅咒者的死亡，以解除法律的诅咒，自愿作为祭品献给天主圣父，好借着一份自愿的祭品，除去因停止常献祭品而伴随的诅咒。此外，圣咏另一处也提及此祭献：\BibleQuote{牺牲与素祭你不喜欢，却为我预备了一个身体。} 意思是，向拒绝法律祭献的天主圣父，献上祂所接受的身体作为可悦纳的祭品。关于此祭品，圣宗徒如此说：\BibleQuote{因为祂只此一次，将自己献上，就完成了这事。}\BibleRef{希 7:27} 借着这圣洁、完美的牺牲祭献，为人类获得了完全的救恩。\par

14. 然后，祂为这一切事的完成感谢天主圣父：\BibleQuote{上主，我要称谢祢的名，因为祢的名是美善的，祢拯救我脱离了一切困苦。} 祂使每一句话都严格应验。因此，起初祂曾说：\BibleQuote{天主，求祢以祢的名拯救我。} 在祈祷蒙垂听后，理当相应地感恩，以使借其名祈求拯救的那一位得着称颂，并且正如祂曾求助于对抗那些起来攻击祂的外邦人，如今在喜乐迸发中记载祂已获得帮助：\BibleQuote{祢拯救我脱离了一切困苦。} 接着，关于那些暴力寻求祂灵魂的人没有将天主放在眼前的事实，祂宣告了自己永恒拥有不变的天主性：\BibleQuote{我的眼曾俯视我的仇敌。} 因为天主的独生子并未被死亡切断。确实，为将我们整个本性引入自身，祂接受了死亡，即灵魂与身体的明显分离，甚至下到阴间，偿还人类显然必须偿还的债务：但祂复活了，永远存在，以死亡不能使之暗淡的眼俯视仇敌，被高举到天主的荣耀中，并在成为人子后再次成为圣子，正如祂最初成为人子时已是圣子，凭着祂复活的荣耀。祂俯视那些祂曾对他们说过的仇敌：\BibleQuote{拆毁这座圣殿，三天内我要把它重建起来。}\BibleRef{若 2:19} 因此，如今祂身体的圣殿已重建，祂从高高的宝座上审视那些曾寻求祂灵魂的人，远超人死亡的权势，从天俯视那些造成祂死亡的人——祂，受了死却不能死的天主而人，我们的主耶稣基督，永受赞美，直到永远。阿们。\par

\SourceRecord{Translated by E.W. Watson and L. Pullan.；Nicene and Post-Nicene Fathers, Second Series；Vol. 9.；Edited by Philip Schaff and Henry Wace.；Buffalo, NY: Christian Literature Publishing Co.,；1899.；<http://www.newadvent.org/fathers/3303053.htm>.}

% structure: p0001 source=h1 target=section reason=page-title

\section{《圣咏》第一三一篇讲道词}

\BibleQuote{上主，我的心不傲慢，我的眼不高举。}\par

1. 这篇圣咏很短，需要的是分析性的处理而非讲道式的处理，它教导我们谦卑和温良的功课。既然我们在许多其他地方已经谈论过谦卑，这里就不必重复相同的内容。当然，我们应当牢记，当我们听到先知如此将谦卑等同于最高行为的成就时，我们的信德是何等需要谦卑：\BibleQuote{上主，我的心不傲慢。}因为一颗忧苦的心在天主眼中是最尊贵的祭献。因此，心不可因顺遂而高举，却要因敬畏天主而谦卑地保持在温良的界限内。\par

2. \Quote{我的眼也不高举。}希腊文在这里的严格含义传达出不同的意思：\OriginalTerm{我的眼也没有被高举}{\GreekText{οὐδὲ} \GreekText{ἐμετεωρίσθησαν} \GreekText{οἱ} \GreekText{ὀφθαλμοί} \GreekText{μου}}，即没有被高举从一个对象去看另一个。然而，眼必须高举，以顺从先知的话：\BibleQuote{举起你们的眼，看谁创造了这一切。}\BibleRef{依 40:26} 主在福音中说：\BibleQuote{举起你们的眼，观看田地，已经发白，可以收割了。}\BibleRef{若 4:35} 因此，眼是要高举的，但不是要转移视线去看别处，而是一旦被举起，就始终凝视那对象。\par

3. 接着是：\Quote{我没有行大事，也没有行超过我的奇事。}行卑贱之事，而不停留于奇事，是最危险的。天主的言语是伟大的；祂本身是至高奇妙的：那么圣咏作者怎能以不行大事和奇事为善功而自夸呢？正是加上了\Quote{超过我的}这几个字，表明他所行的不是世人通常视为伟大和奇妙的事。因为达味，既是先知又是君王，曾经卑微、被轻视、不配坐在他父亲的席上；但他蒙天主恩宠，受傅为君王，被感动而说预言。他的王国没有使他傲慢，他没有被仇恨所动：他爱那些迫害他的人，他尊敬死去的敌人，他宽恕了他那乱伦和杀人的子女。作为君王，他被轻视；作为父亲，他受伤；作为先知，他受折磨；然而他没有像先知那样呼求报复，没有像父亲那样施加惩罚，没有像君王那样报复侮辱。因此，他没有行超过他的大事和奇事。\par

4. 让我们看看接下来的内容：\Quote{若我不谦卑自持，却高举了我的灵魂。}先知是多么矛盾！他不高举他的心，却高举他的灵魂。他没有行超过他的大事和奇事；但他的思想却不卑微。他思想崇高，心却谦卑。他在自己的事务上是谦卑的，但他的思想并不谦卑。因为他的思想达于天上，他的灵魂被高举至高处。但他的心——按福音所说，\BibleQuote{由心里发出来的是恶念、凶杀、奸淫、淫乱、偷窃、妄证、毁谤}\BibleRef{玛 15:19}——却是谦卑的，被压下在温良的柔和轭下。因此，我们必须在谦卑与高举之间走一条中间路线，使我们心谦卑，但灵魂和思想却高举。\par

5. 然后他接着说：\BibleQuote{如同断乳的幼儿在他母亲的怀中，你同样会酬报我的灵魂。}我们得知，当依撒格断乳时，亚巴郎摆设了宴席，因为他既已断乳，就接近少年时期，超越了奶食。宗徒用知识的奶喂养所有信德不完善、在天主的事上仍为婴孩的人。因此，不再需要奶标志着最大的进步。亚巴郎以欢乐的宴席宣告他的儿子已可吃较硬的食物，宗徒拒绝将饼给属肉体和在基督里为婴孩的人。因此，先知祈求天主，因为他没有高举他的心，也没有行超过他的大事和奇事，因为他不曾谦卑自持，却高举了他的灵魂，求天主酬报他的灵魂，如同躺在母亲怀中的断乳幼儿：也就是说，他因着已记载的作为，如今已超越了奶的阶段，从而堪受完美、属天、永生的食粮的赏报。\par

6. 但他不仅为自己要求这来自天上的永生食粮，他也鼓励全人类盼望它，说：\BibleQuote{从今直到永远，以色列要仰赖上主。}他对我们的盼望没有设定期限，他命令我们信靠的期待伸展至无限。我们要永远盼望，借着我们在现今生命中的盼望——这盼望在基督耶稣我们的主内——赢得未来生命的盼望。祂是当受赞美的，直到永远。阿们。\par

\SourceRecord{Translated by E.W. Watson and L. Pullan.；Nicene and Post-Nicene Fathers, Second Series；Vol. 9.；Edited by Philip Schaff and Henry Wace.；Buffalo, NY: Christian Literature Publishing Co.,；1899.；<http://www.newadvent.org/fathers/3303131.htm>.}
